\begingroup
\raggedright
\GaussContentsEntry{\SourceLeaf{666}\noindent\hangindent=1em [\textsc{Grahm}'s planet observed]\GaussContentsLeaders{\mbox{.\ .\ .\ .\ .\ .\ .\ .\ .\ .\ .\ .\ .\ .\ .\ .\ .}}\GaussContentsReference{\mbox{1848 June 10}\quad\mbox{Page 477}}}
\GaussContentsEntry{\noindent\hangindent=1em [\textsc{Graham}'s planet observed. Elements by \textsc{Goldschmidt}]\GaussContentsLeaders{\mbox{.\ .\ .\ .\ .\ .}}\GaussContentsReference{\mbox{1848 June 29}\quad\mbox{--- 477}}}
\GaussContentsEntry{\noindent\hangindent=1em [\textsc{Hind}'s star observed]\GaussContentsLeaders{\mbox{.\ .\ .\ .\ .\ .\ .\ .\ .\ .\ .\ .\ .\ .\ .\ .\ .\ .\ .}}\GaussContentsReference{\mbox{1848 Nov. 2}\quad\mbox{--- 478}}}
\GaussContentsEntry{\noindent\hangindent=1em [Parthenope observed]\GaussContentsLeaders{\mbox{.\ .\ .\ .\ .\ .\ .\ .\ .\ .\ .\ .\ .\ .\ .\ .\ .\ .\ .}}\GaussContentsReference{\mbox{1850 June 20}\quad\mbox{--- 478}}}
\GaussContentsEntry{\noindent\hangindent=1em [Victoria observed]\GaussContentsLeaders{\mbox{.\ .\ .\ .\ .\ .\ .\ .\ .\ .\ .\ .\ .\ .\ .\ .\ .\ .\ .\ .}}\GaussContentsReference{\mbox{1850 Nov. 22}\quad\mbox{--- 478}}}
\GaussContentsEntry{\noindent\hangindent=1em [Solar eclipse observed]\GaussContentsLeaders{\mbox{.\ .\ .\ .\ .\ .\ .\ .\ .\ .\ .\ .\ .\ .\ .\ .\ .}}\GaussContentsReference{\mbox{1851 Aug. 28}\quad\mbox{--- 479}}}
\GaussContentsEntry{\noindent\hangindent=1em [Psyche]\GaussContentsLeaders{\mbox{.\ .\ .\ .\ .\ .\ .\ .\ .\ .\ .\ .\ .\ .\ .\ .\ .\ .\ .\ .\ .\ .\ .\ .\ .}}\GaussContentsReference{\mbox{1852 July 27}\quad\mbox{--- 479}}}
\bigskip\Needspace{5\baselineskip}\noindent\textit{Reviews and notices of works by other authors.}\par\medskip
\GaussContentsEntry{\noindent\hangindent=1em \textsc{Rohde}: On the Masses of the Planets\GaussContentsLeaders{\mbox{.\ .\ .\ .\ .\ .\ .\ .\ .\ .\ .\ .\ .\ .}}\GaussContentsReference{\mbox{1805 October}\quad\mbox{--- 483}}}
\GaussContentsEntry{\noindent\hangindent=1em On \textsc{La Place}'s expression for measuring heights with a barometer\GaussContentsLeaders{\mbox{.\ .\ .}}\GaussContentsReference{\mbox{1806 May 18}\quad\mbox{--- 485}}}
\GaussContentsEntry{\noindent\hangindent=1em Posthumous papers: von \textsc{Lindenau} to \textsc{Gauss}\GaussContentsLeaders{\mbox{.\ .\ .\ .\ .\ .\ .\ .\ .\ .\ .\ .\ .\ .\ .}}\GaussContentsReference{\mbox{1806 Feb. 17}\quad\mbox{--- 486}}}
\GaussContentsEntry{\noindent\hangindent=1em Posthumous papers: von \textsc{Zach} to \textsc{Gauss}\GaussContentsLeaders{\mbox{.\ .\ .\ .\ .\ .\ .\ .\ .\ .\ .\ .\ .\hspace{.6em}\ .\ .\ .}}\GaussContentsReference{\mbox{1807 April 6}\quad\mbox{--- 486}}}
\GaussContentsEntry{\noindent\hangindent=1em Extract from a letter [On \textsc{Güssmann}: Comet Orbits]\GaussContentsLeaders{\mbox{.\ .\ .\ .}}\GaussContentsReference{\mbox{1807 March 10}\quad\mbox{--- 486}}}
\GaussContentsEntry{\noindent\hangindent=1em \textsc{De Zach}: Tabulae speciales aberrationis et nutationis in ascensionem rectam et in declinationem ad supputandas stellarum fixarum positiones, sive apparentes, sive veras, una cum insigniorum 494 stellarum zodiacalium catalogo novo in specula astronomica Ernestina ad initium anni 1800 constructo, cum aliis tabulis eo spectantibus\GaussContentsLeaders{\mbox{.\ .\ .\ .\ .\ .\ .\ .\ .\ .}}\GaussContentsReference{\mbox{1808 Jan. 2}\quad\mbox{--- 489}}}
\GaussContentsEntry{\noindent\hangindent=1em Connoissance des tems, ou des mouvements célestes à l'usage des astronomes et des navigateurs pour l'an 1808 publiée par le bureau des longitudes. Paris 1806. Pour l'an 1809. Paris 1807\GaussContentsLeaders{\mbox{.\ .\ .\ .\ .\ .\ .\ .\ .\ .}}\GaussContentsReference{\mbox{1808 April 7}\quad\mbox{--- 492}}}
\GaussContentsEntry{\noindent\hangindent=1em Antoine Cagnoli: Catalogue de 501 étoiles suivi des tables relatives d'Aberration et de Nutation. Modena 1807\GaussContentsLeaders{\mbox{.\ .\ .\ .\ .\ .\ .\ .\ .\ .\ .\ .}}\GaussContentsReference{\mbox{1808 May 14}\quad\mbox{--- 496}}}
\GaussContentsEntry{\noindent\hangindent=1em \textsc{Pasquich}: Account of my proposals for advancing astronomy at the Royal University Observatory at Ofen\GaussContentsLeaders{\mbox{.\ .\ .\ .\ .\ .}}\GaussContentsReference{\mbox{1808 May 28}\quad\mbox{--- 497}}}
\GaussContentsEntry{\noindent\hangindent=1em \textsc{von} \textsc{Lindenau}: On the Present State of the Seeberg Observatory\GaussContentsLeaders{\mbox{.}}\GaussContentsReference{\mbox{1808 July 2}\quad\mbox{--- 498}}}
\GaussContentsEntry{\noindent\hangindent=1em \textsc{Laplace}: Exposition du système du monde. Troisième édition. 1808\GaussContentsLeaders{\mbox{.\ .\ .}}\GaussContentsReference{\mbox{1808 July 23}\quad\mbox{--- 500}}}
\GaussContentsEntry{\noindent\hangindent=1em \textsc{Bode}: Astronomical Yearbook for the Year 1810. Berlin 1807\GaussContentsLeaders{\mbox{.\ .\ .\ .\ .}}\GaussContentsReference{\mbox{1808 Sept. 26}\quad\mbox{--- 504}}}
\GaussContentsEntry{\noindent\hangindent=1em \textsc{Picot}: Sur les vingt-une dernières comètes et les nouvelles planètes. Bibliothèque britannique 1808 May. Geneva\GaussContentsLeaders{\mbox{.\ .\ .\ .\ .\ .\ .\ .\ .\ .\ .\ .}}\GaussContentsReference{\mbox{1808 Oct. 13}\quad\mbox{--- 507}}}
\GaussContentsEntry{\noindent\hangindent=1em \textsc{Bouvard}: Tables astronomiques publiées par le bureau des longitudes de France. Nouvelles tables de Jupiter et de Saturne calculées d'après la théorie de M. \textsc{Laplace}, et suivant la division decimale de l'angle droit. Paris 1808\GaussContentsLeaders{\mbox{.\ .\ .\ .\ .\ .\ .\ .\ .\ .\ .\ .\ .\ .\ .\ .\ .\ .\ .\ .}}\GaussContentsReference{\mbox{1808 Nov. 5}\quad\mbox{--- 508}}}
\GaussContentsEntry{\noindent\hangindent=1em \textsc{Jabbo Oltmans}: Conspectus longitudinum et latitudinum geographicarum, per decursum annorum 1799 ad 1804 in plaga aequinoctiali ab \textsc{Alex. de Humboldt} astronomice observatarum, calculo subjecit. Paris and Tübingen \GaussContentsReference{\mbox{1808 Dec. 3}\quad\mbox{--- 509}}}
\GaussContentsEntry{\noindent\hangindent=1em \textsc{Harding}: New Star Atlas. First Instalment. 1809. Hamburg, from Perthes \GaussContentsReference{\mbox{1809 April 17}\quad\mbox{--- 510}}}
\GaussContentsEntry{\noindent\hangindent=1em Voyage \textsc{de Humboldt} et \textsc{Bonpland}. Quatrième Partie, Astronomie et Magnetisme. Premier volume, contenant un recueil d'observations astronomiques, d'opérations trigonométriques, et de mésures barométriques, faites pendant le cours d'un voyage aux régions équinoxiales du Nouveau-Continent, depuis 1799 jusqu'en 1803. Instalment I, 1808. Instalment II, 1809. Paris and Tübingen\GaussContentsLeaders{\mbox{.\ .\ .\ .\ .\ .\ .\ .\ .\ .\ .\ .\ .\ .}}\GaussContentsReference{\mbox{1809 June 10}\quad\mbox{--- 512}}}
\endgroup
